% ============================================================
%  Sobre o Autor
% ============================================================
\chapter*{Sobre o Autor}
\addcontentsline{toc}{chapter}{Sobre o Autor}
\thispagestyle{plain}

\noindent
\textbf{Wilian Fran\c{c}a Costa} \'e brasileiro, natural de S\~ao Paulo.
Bacharel em Ci\^encia da Computa\c{c}\~ao pela Unesp (2006), Mestre em
Ci\^encias pelo Programa de P\'os-Gradua\c{c}\~ao em Engenharia Eletr\^onica
e Computa\c{c}\~ao do Instituto Tecnol\'ogico de Aeron\'autica --- ITA (2010)
e Doutor em Ci\^encias no Programa de Engenharia El\'etrica ---
\'area de concentra\c{c}\~ao Engenharia de Computa\c{c}\~ao --- da Escola
Polit\'ecnica da Universidade de S\~ao Paulo --- POLI-USP (2016).

Iniciou a vida profissional como t\'ecnico em eletr\^onica e eletricidade,
tendo concludo o t\'ecnico em Processamento de Dados pelo Col\'egio
S\~ao Lu\'is -- Jesu\'itas (2000) antes de ingressar na gradua\c{c}\~ao.
Essa trajet\'oria marcada pela pr\'atica industrial fundamentou sua
vis\~ao de que tecnologia deve ser ao mesmo tempo rigorosa e aplic\'avel.

Ao longo de sua carreira atuou como professor universit\'ario em
disciplinas de Ci\^encia da Computa\c{c}\~ao, Engenharia de
Computa\c{c}\~ao e Sistemas de Informa\c{c}\~ao. Participou de projetos
de pesquisa envolvendo geoprocessamento, navega\c{c}\~ao aut\^onoma e
mapeamento simult\^aneo (SLAM), agricultura de precis\~ao e
arquiteturas orientadas a servi\c{c}os. Sua disserta\c{c}\~ao de
mestrado no ITA gerou uma publica\c{c}\~ao no \emph{Latin American
Robotics Symposium and Intelligent Robotics Meeting} (LARS~2010), e sua
tese de doutorado na USP abordou m\'etodos geoespaciais aplicados \`a
agricultura de precis\~ao.

Atualmente \'e Analista S\^enior de Ci\^encia de Dados no Ita\'u
Unibanco, onde pesquisa e implementa sistemas baseados em \emph{Large
Language Models} (LLMs), agentes aut\^onomos e arquiteturas de
recupera\c{c}\~ao aumentada por gera\c{c}\~ao (RAG) em ambientes
corporativos regulados.  Sua trajet\'oria re\'une forma\c{c}\~ao
acad\^emica de alto n\'ivel, experi\^encia industrial profunda e
produ\c{c}\~ao cient\'ifica em confer\^encias e peri\'odicos nacionais
e internacionais.

Este livro \'e resultado direto dessa jornada: cada cap\'itulo nasce
de decis\~oes t\'ecnicas reais, implementa\c{c}\~oes testadas em
produ\c{c}\~ao e licoes aprendidas no cruzamento entre pesquisa e
engenharia.

\bigskip
\noindent
\begin{tabular}{@{}ll}
\textbf{Forma\c{c}\~ao}  & Doutor em Engenharia de Computa\c{c}\~ao --- POLI-USP (2016) \\
                          & Mestre em Engenharia Eletr\^onica e Computa\c{c}\~ao --- ITA (2010) \\
                          & Bacharel em Ci\^encia da Computa\c{c}\~ao --- Unesp (2006) \\[0.4em]
\textbf{Atua\c{c}\~ao}   & Analista S\^enior de Ci\^encia de Dados --- Ita\'u Unibanco \\
                          & Pesquisador em LLMs, RAG e Agentes Aut\^onomos \\
\end{tabular}
